\chapter{Evren, Örneklem, Parametre, İstatistik ve Veri Türleri}
\label{ch:population-data}

\begin{prismbox}[PrismLight]{Öğrenme çıktıları}
Bu bölümün sonunda hedef evren ile erişilebilir evreni, parametre ile
istatistiği ve kategorik veri ile nicel veriyi birbirinden ayırabilecek;
araştırma sorusunu uygun analiz birimi ve değişkenlerle ifade edebileceksiniz.
\end{prismbox}

\begin{prerequisitebox}
Bölüm 1'deki araştırma sorusu, gözlem birimi ve değişken kavramlarının
bilinmesi beklenir. Oran ve ortalama kavramlarına giriş düzeyinde aşinalık
yeterlidir.
\end{prerequisitebox}

Araştırma sorusunun değişkenlere dönüştürülmesi ve veri üretim zinciri için
Bölüm~\ref{ch:research}, özellikle Bölüm~\ref{sec:research-cycle} temel alınır.

\section{Araştırmanın birimi ve hedef evren}
\label{sec:population-target}

İstatistiksel bir araştırma önce \emph{kimin} veya \emph{neyin} incelendiğini
belirler. Üzerinde ölçüm yapılan her öğrenci, sınıf, okul ya da deney birimi
bir \textbf{gözlem birimi}dir. Sonucun genellenmek istendiği bütün birimler
\textbf{hedef evreni}, araştırmacının fiilen ulaşabildiği bölüm ise
\textbf{erişilebilir evreni} oluşturur.

Örneğin Türkiye'deki bütün ortaokul öğrencileri hedeflenirken veri yalnız bir
ildeki devlet okullarından toplanmışsa erişilebilir evren hedef evrenden daha
dardır. Bu ayrım, sonuçların hangi topluluğa genellenebileceğini belirler.

\section{Evren ve örneklem}

Evrenin tamamına ilişkin ölçüm yapılmasına \textbf{tam sayım} denir. Zaman,
maliyet ve erişim sınırlılıkları nedeniyle çoğu araştırmada evrenden bir
\textbf{örneklem} seçilir. Örneklemin amacı yalnız daha az birim incelemek
değil, seçilen birimlerden hareketle hedef evren hakkında ölçülebilir
belirsizlik altında bilgi üretmektir.
Örnekleme çerçevesi, seçilme olasılığı ve tasarımın genellemeye etkisi için
temel başvuru kaynakları arasında \citet{cochran1977} ile \citet{lohr2021}
yer alır.

\begin{mistakebox}
Örneklemin evrene benzemesi yalnız örneklem büyüklüğüne bağlı değildir.
Seçim mekanizması bazı grupları sistematik olarak dışlıyorsa çok büyük bir
örneklem de yanıltıcı olabilir.
\end{mistakebox}

Şekil~\ref{fig:population-nesting}'de örneklem, erişilebilir evrenin;
erişilebilir evren ise hedef evrenin içinde gösterilir. Örnekte seçilen
öğrenci sayısını artırmak, kapsanmayan illeri veya özel okulları kendiliğinden
temsil edilir hâle getirmez. Alanların büyüklüğü öğrenci sayısıyla orantılı
değildir.

\begin{figure}[H]
\centering
\begin{tikzpicture}[font=\small]
  \draw[PrismBlue,fill=PrismLight,rounded corners=2mm]
    (-5.25,0) rectangle (5.25,4.9);
  \node[align=center] at (0,4.35)
    {\textbf{Hedef evren}\\Türkiye'deki bütün ortaokul öğrencileri};
  \draw[PrismBlue,dashed,fill=white,rounded corners=2mm]
    (-4.6,0.35) rectangle (4.6,3.7);
  \node[align=center] at (0,3.1)
    {\textbf{Erişilebilir evren}\\Bir ildeki devlet ortaokullarının öğrencileri};
  \draw[PrismBlue,line width=1pt,fill=PrismNoteBlue,rounded corners=2mm]
    (-3.7,0.75) rectangle (3.7,2.45);
  \node[align=center] at (0,1.9)
    {\textbf{Örneklem}\\Bu okullardan seçilen öğrenciler};
  \foreach \x in {-2.4,-1.6,-0.8,0,0.8,1.6,2.4}
    \fill[PrismBlue] (\x,1.12) circle (2.3pt);
\end{tikzpicture}
\caption{Hedef evren, erişilebilir evren ve örneklem: bir kapsam sınırlılığı
örneği. Noktalar seçilen öğrencileri simgeler.}
\label{fig:population-nesting}
\end{figure}

\section{Parametre ve istatistik}
\label{sec:population-parameter}

Evrene ait sayısal özellik \textbf{parametre}, örneklemden hesaplanan karşılığı
ise \textbf{istatistik} olarak adlandırılır. Parametre sabit fakat çoğu zaman
bilinmeyendir; istatistik seçilen örnekleme göre değişir.

\begin{table}[H]
\centering
\caption{Evren büyüklükleri ve örneklem karşılıkları.}
\label{tab:population-sample-examples}
\begin{tabular}{lll}
\toprule
Araştırma büyüklüğü & Parametre & İstatistik \\
\midrule
Ortalama & $\mu$ & $\bar x$ \\
Oran & $p$ & $\hat p$ \\
Varyans & $\sigma^2$ & $s^2$ \\
Korelasyon & $\rho$ & $r$ \\
\bottomrule
\end{tabular}
\end{table}

Bir fakültedeki bütün öğrencilerin haftalık çalışma süresi ortalaması $\mu$
ile gösterilen bir parametredir. Rastgele seçilen 120 öğrencinin ortalaması
$\bar x=6{,}4$ saat ise bu değer bir istatistiktir ve $\mu$ hakkında tahmin
üretmek için kullanılır.

\section{Veri türleri}
\label{sec:population-data-types}

\textbf{Kategorik değişkenler} birimleri sınıflara ayırır. Nominal değişkende
kategorilerin doğal sırası yoktur; sıralı değişkende ise sıra vardır fakat
kategoriler arasındaki uzaklık zorunlu olarak eşit değildir.
\textbf{Nicel değişkenler} sayısal miktar bildirir. Sayım sonucu elde edilen
değerler genellikle kesikli, ölçüm sonucu elde edilenler sürekli kabul edilir.

\begin{table}[H]
\centering
\caption{Veri türleri için eğitim araştırması örnekleri.}
\label{tab:data-type-examples}
\begin{tabularx}{\textwidth}{@{}p{0.24\textwidth}XX@{}}
\toprule
Tür & Örnek & Uygun ilk özet \\
\midrule
Nominal & Okul türü & Frekans ve yüzde \\
Sıralı & Memnuniyet düzeyi & Frekans, medyan, yüzdelik \\
Kesikli nicel & Doğru yanıt sayısı & Ortalama, medyan, dağılım \\
Sürekli nicel & Tamamlama süresi & Ortalama, standart sapma, histogram \\
\bottomrule
\end{tabularx}
\end{table}

Sayısal kod kullanılması bir değişkeni kendiliğinden nicel yapmaz. Okul
türlerinin 1, 2 ve 3 ile kodlanması yalnız veri girişini kolaylaştırır; bu
kodların ortalaması anlamlı bir okul türü ölçüsü değildir.

\section{Örnekleme yöntemlerine ilk bakış}
\label{sec:population-sampling-intro}

Örneklem yalnız kaç birimden oluştuğuyla değil, birimlerin nasıl seçildiğiyle
de tanımlanır. \textbf{Basit rastgele örneklemede} her birimin bilinen bir
seçilme olasılığı vardır. \textbf{Tabakalı örneklemede} fakülte, sınıf düzeyi
veya okul türü gibi önemli alt gruplardan ayrı seçim yapılarak temsil
güçlendirilir. \textbf{Küme örneklemesinde} bireyler yerine sınıf veya okul
gibi doğal gruplar seçilir. \textbf{Kolayda} ve \textbf{gönüllü}
örneklemelerde ise ulaşılabilen ya da katılmayı seçen bireyler kullanılır.

Olasılıklı yöntemler örnekleme belirsizliğinin hesaplanmasına ve hedef evrene
genelleme yapılmasına daha güçlü bir temel sağlar. Olasılıksız yöntemler bazı
araştırma koşullarında gerekli olabilir; ancak genelleme sınırları açıkça
belirtilmelidir. Bu yöntemlerin tasarım ayrıntıları Bölüm 6'da ele alınacaktır.

\section{Temsil, kapsam ve yanlılık}
\label{sec:population-coverage}

Örnekleme çerçevesi, birimlerin seçildiği fiilî listedir. Çerçevenin hedef
evreni eksik kapsaması \textbf{kapsam yanlılığına}; seçilen kişilerin önemli
bir bölümünün yanıt vermemesi \textbf{yanıtlamama yanlılığına}; yalnız
isteklilerin katılması ise \textbf{gönüllülük yanlılığına} yol açabilir.

Örneğin bütün üniversite öğrencileri hedeflenirken anketin yalnız mezuniyet
aşamasındaki öğrencilere e-posta ile gönderilmesi, örneklem büyüklüğü yüksek
olsa bile alt sınıfları dışarıda bırakır. İstatistiksel belirsizlik ile
sistematik yanlılık ayrı sorunlardır: daha büyük örneklem rastgele
dalgalanmayı azaltabilir, fakat kapsam hatasını kendiliğinden düzeltemez.

\begin{researchbox}
Bir sonuç raporlanırken örneklem hacminin yanında örnekleme çerçevesi, seçim
yöntemi, katılım oranı ve hedef evrenden olası sapmalar da belirtilmelidir.
``Katılımcılar'' ile ``bütün öğrenciler'' aynı topluluk değildir.
\end{researchbox}

\section{Çözümlü sınıflandırma örneği}

Bir araştırmada öğrencinin programı, sınıf düzeyi, sınav puanı ve sınavı
tamamlama süresi kaydedilsin. Program nominal kategorik, sınıf düzeyi sıralı
kategorik, sınav puanı kesikli nicel ve süre sürekli nicel değişkendir. Bütün
öğrencilerin ortalama sınav puanı parametre; seçilen öğrencilerin ortalaması
istatistiktir.

\section{Yazılım uygulaması: kategorik ve nicel değişkenler}

\begin{softwarebox}
\begin{lstlisting}
import pandas as pd

data = pd.DataFrame({
    "okul_turu": ["Devlet", "Ozel", "Devlet", "Devlet"],
    "sinif": [1, 2, 1, 3],
    "puan": [72, 81, 68, 77]
})
data["okul_turu"] = data["okul_turu"].astype("category")
data["sinif"] = pd.Categorical(
    data["sinif"], categories=[1, 2, 3], ordered=True)

print(data.dtypes)
print(data["okul_turu"].value_counts(normalize=True))
print(data["puan"].mean())
\end{lstlisting}
Kodlama, okul türü ile puanı aynı sayısal işlem türüne zorlamak yerine her
değişkenin yapısını korur.
\end{softwarebox}

\begin{assumptionbox}
Veri türü yalnız dosyadaki teknik biçimden belirlenmez. Bir değişken yazılımda
sayısal görünse bile araştırmadaki anlamı kimlik, kategori veya sıralama
olabilir. Analiz kararı kavramsal tanıma dayanmalıdır.
\end{assumptionbox}

\section{Çözümlü problemler}

\subsection*{Problem 1: Hedef ve erişilebilir evren}

Bir araştırmacı Türkiye'deki lise öğrencilerinin çevrim içi öğrenmeye yönelik
tutumlarını incelemek istemekte, ancak veriyi yalnız İstanbul'daki devlet
liselerinden toplamaktadır. Hedef evreni, erişilebilir evreni, örnekleme
çerçevesini ve güvenli genelleme sınırını belirleyin.

\textbf{Çözüm.} Araştırma sorusunda belirtilen hedef evren Türkiye'deki bütün
lise öğrencileridir. Araştırmacının fiilen ulaşabildiği evren İstanbul'daki
devlet liselerinde öğrenim gören öğrencilerdir. Öğrencilerin seçildiği güncel
okul veya öğrenci listesi örnekleme çerçevesidir. Özel lise öğrencileri,
diğer illerdeki öğrenciler ve okul dışında kalan gençler bu çerçevede
bulunmamaktadır.

Bu nedenle sonuçların doğrudan Türkiye'deki bütün lise öğrencilerine
genellenmesi desteklenmez. Seçim İstanbul'daki devlet liselerini temsil edecek
biçimde yapılmışsa en güvenli çıkarım bu toplulukla sınırlıdır. Daha geniş
genelleme için bölge, okul türü ve sosyoekonomik yapı bakımından daha kapsamlı
bir örnekleme çerçevesi gerekir.

\textbf{Raporlama örneği.} ``Örneklem İstanbul'daki devlet liselerinden
seçilmiştir; bu nedenle bulguların özel liselere ve diğer illerdeki öğrencilere
genellenmesi sınırlıdır.''

\subsection*{Problem 2: Gözlem birimi ile analiz birimi}

Yirmi okuldan seçilen 60 sınıfta toplam 1.800 öğrencinin matematik puanı
ölçülmüştür. Bir satır bir öğrenciyi göstermektedir. Gözlem birimi ve analiz
birimi nedir? Öğrencilerin tamamen bağımsız kabul edilmesi neden sorunlu
olabilir?

\textbf{Çözüm.} Ölçüm doğrudan öğrenciler üzerinde yapıldığı için gözlem
birimi öğrencidir. Öğrenci puanları karşılaştırılıyorsa temel analiz birimi de
öğrencidir. Bununla birlikte öğrenciler sınıfların, sınıflar da okulların
içinde kümelenmiştir. Aynı sınıftaki öğrenciler ortak öğretmen, ders ortamı ve
akran etkileri nedeniyle farklı sınıflardaki öğrencilerden daha çok
birbirlerine benzeyebilir.

Bu yapı gözlem sayısını değiştirmez; veri setinde yine 1.800 öğrenci vardır.
Ancak 1.800 tamamen bağımsız öğrenciden elde edilecek bilgiyle aynı miktarda
bilgi bulunduğu varsayımı aşırı iyimser standart hatalara yol açabilir.
Analizde sınıf ve okul kümelenmesini dikkate alan yöntemler gerekebilir.

\textbf{Sık hata.} Veri dosyasındaki satır sayısını doğrudan bağımsız örneklem
hacmi kabul etmek, çok düzeyli veya tekrarlı ölçümlü yapılarda yanıltıcıdır.

\subsection*{Problem 3: Parametre, istatistik ve tahmin hedefi}

Bir ilçedeki 8. sınıf öğrencilerinin liseye geçiş sınavında başarı oranı
araştırılmaktadır. İlçede 4.200 öğrenci vardır. Rastgele seçilen 300 öğrencinin
198'i belirlenen başarı ölçütünü karşılamıştır. Parametreyi, istatistiği ve
örneklem oranını bulun.

\textbf{Çözüm.} İlçedeki 4.200 öğrencinin başarı ölçütünü karşılama oranı
bilinmeyen evren parametresi $p$'dir. Örneklemde başarılı olan öğrenci sayısı
198, örneklem hacmi 300'dür. Örneklem oranı
\[
\hat p=\frac{198}{300}=0{,}66
\]
olur. Dolayısıyla öğrencilerin yüzde 66'sı örneklemde başarı ölçütünü
karşılamıştır. Bu yüzde, evrendeki gerçek oran $p$ için bir tahmindir; $p$'nin
kesin olarak yüzde 66 olduğunu göstermez.

\textbf{Raporlama örneği.} ``Rastgele seçilen 300 öğrencinin 198'i başarı
ölçütünü karşılamış, örneklem başarı oranı yüzde 66 bulunmuştur. Bu değer
ilçedeki 8. sınıf öğrencilerinin başarı oranı için bir nokta tahminidir.''

\subsection*{Problem 4: Değişken türü ve veri sözlüğü}

Bir araştırma dosyasında şu sütunlar bulunmaktadır: öğrenci kodu, okul türü
(1=devlet, 2=özel), devamsızlık günü, sınıf başarı sırası ve sınavı tamamlama
süresi. Her değişkeni sınıflandırın ve uygun bir veri doğrulama kuralı yazın.

\textbf{Çözüm.}

\begin{center}
\small
\begin{tabularx}{\textwidth}{@{}p{0.22\textwidth}p{0.23\textwidth}X@{}}
\toprule
Değişken & Tür & Örnek doğrulama kuralı \\
\midrule
Öğrenci kodu & Nominal kimlik & Her öğrenci için benzersiz olmalı \\
Okul türü & Nominal kategorik & Yalnız 1 veya 2 değerini almalı \\
Devamsızlık günü & Kesikli nicel & Negatif olmamalı; dönem gününü aşmamalı \\
Başarı sırası & Sıralı & 1 ile öğrenci sayısı arasında olmalı \\
Tamamlama süresi & Sürekli nicel & Pozitif ve sınav süresiyle uyumlu olmalı \\
\bottomrule
\end{tabularx}
\end{center}

Öğrenci kodunun ve okul türünün rakamla yazılması bunları nicel değişken
yapmaz. Kod ortalaması veya okul türü kodlarının standart sapması anlamlı bir
araştırma özeti değildir. Devamsızlık ve süre için merkezi eğilim ile yayılım
özetleri kullanılabilir; fakat önce olanaksız değerler araştırılmalıdır.

\subsection*{Problem 5: Örnekleme yöntemi seçimi}

Bir üniversite, öğrencilerin kütüphane kullanım süresini tahmin etmek
istemektedir. Öğrencilerin yüzde 60'ı lisans, yüzde 25'i yüksek lisans ve
yüzde 15'i doktora öğrencisidir. Toplam 400 öğrenci seçilecektir. Her öğrenim
düzeyinin temsili isteniyorsa hangi yöntem uygundur ve orantılı seçimde her
gruptan kaç öğrenci alınmalıdır?

\textbf{Çözüm.} Öğrenim düzeylerinin ayrı ayrı temsil edilmesi istendiği için
tabakalı örnekleme uygundur. Öğrenim düzeyi tabaka değişkeni olarak kullanılır.
Orantılı dağılımda
\[
400(0{,}60)=240,\qquad
400(0{,}25)=100,\qquad
400(0{,}15)=60
\]
olduğundan 240 lisans, 100 yüksek lisans ve 60 doktora öğrencisi seçilir.
Her tabakanın içinde öğrenciler rastgele seçilmelidir.

Doktora öğrencileri için ayrıca ayrıntılı analiz yapılacaksa orantısız biçimde
daha fazla doktora öğrencisi seçilebilir. Bu durumda üniversitenin tamamına
ilişkin bir ortalama hesaplanırken farklı seçilme olasılıklarını yansıtan
ağırlıklar gerekebilir.

\begin{outputguidebox}
Bir veri çıktısını yorumlamadan önce değişken etiketleri, kategori kodları,
eksik değer tanımları ve analiz birimi kontrol edilmelidir. ``1'' değerinin
bir puan mı, kategori mi yoksa eksik veri kodu mu olduğu bilinmeden doğru
istatistik üretilemez.
\end{outputguidebox}

\section{Literatür verisiyle yazılım uygulamaları}
\label{sec:spss-b02}

\textbf{Amaç ve veri.} \texttt{b02.csv} ile aynı 395 öğrencinin okul ve
çalışma süresi kategorilerini inceleyin. Bu, \citet{cortez2008data} verisinin
bir başka uygulamasıdır; yeni bir bağımsız örneklem değildir.


\textbf{Ortak veri.} Üç yazılımda aynı kayıtlar ve değişkenler kullanılır.
\nolinkurl{companion/spss/csv/b02.csv} ile B01 pilot paketindeki
\texttt{b01.csv} dosyaları bayt düzeyinde aynıdır. Paylaşılan B02
uygulamasında R ve Python mevcut \texttt{b01.csv} dosyasını, SPSS ise
B01 CSV içe aktarımında kaydedilen \texttt{b01.sav} dosyasını kullanmıştır.
Bu karşılaştırma Excel içe aktarmanın doğrulandığı anlamına gelmez.
Kodları \texttt{companion} klasörünü içeren paket kökünden çalıştırın;
veri sözlüğü ve hazırlık için bkz. sayfa~\pageref{sec:spss-guide}.

\subsection{IBM SPSS uygulaması}
\textbf{Başlamadan önce.} Doğrulanan CSV yolunu izlemek için B01
verisini açan \texttt{open-b01.sps} dosyasını çalıştırın; bu işlem
\texttt{b01.sav} dosyasını yeniden oluşturur. Aşağıdaki tam blok da
aynı hazırlığı yapıp etiketli dosyayı açar; menü yolu alternatif uygulamadır.

\begin{spssadimlar}
\spssadim{\emph{Variable View} içinde \texttt{studytime} değer etiketlerini inceleyin: 1--4 kodları doğrudan saat sayısı değildir.}
\spssadim{\emph{Analyze $\to$ Descriptive Statistics $\to$ Frequencies} yolunu açın. \texttt{school}, \texttt{sex}, \texttt{studytime} değişkenlerini seçin.}
\spssadim{\emph{Display frequency tables} işaretli kalsın; \emph{OK} seçin. \emph{Frequency}, \emph{Percent} ve \emph{Valid Percent} sütunlarını okuyun.}
\spssadim{\emph{Analyze $\to$ Descriptive Statistics $\to$ Descriptives} yolunda \texttt{age} ve \texttt{g3} değişkenlerini seçin.}
\spssadim{\emph{Options} altında ortalama, standart sapma, en küçük ve en büyük değerleri seçip \emph{Continue $\to$ OK} ile çıktıyı alın. Nominal kodların ortalamasını istemeyin.}
\end{spssadimlar}

{\raggedright
\noindent\textbf{Komutların görevi.} \texttt{FREQUENCIES} kategorik değişkenleri sayar; \texttt{DESCRIPTIVES} yalnız yaş ve notu özetler. \texttt{DISPLAY DICTIONARY} veri tanımlarını kontrol etmeyi sağlar.\par}

\par\noindent\begin{minipage}{\linewidth}
\textbf{Uygulama kodu:} \texttt{b02}. Doğrulanan pilotun veri dosyası:
\texttt{sav/b01.sav} (\texttt{companion/spss} altında).

\begin{lstlisting}[style=prismSPSS]
INSERT FILE='companion/spss/syntax/open-b01.sps'
 ERROR=STOP.
GET FILE='companion/spss/sav/b01.sav'.
FILTER OFF.
WEIGHT OFF.
SPLIT FILE OFF.
DISPLAY DICTIONARY.
FREQUENCIES VARIABLES=school sex studytime.
DESCRIPTIVES VARIABLES=age g3
 /STATISTICS=MEAN STDDEV MIN MAX.
\end{lstlisting}

\end{minipage}\par

\begin{outputguidebox}
\textbf{Paylaşılan SPSS çıktısıyla doğrulanan değerler.}
\texttt{studytime} için 1, 2, 3, 4 kategorilerinin frekansları sırasıyla
105, 198, 65 ve 27'dir. Kategori 2, haftalık 2--5 saat aralığını gösterir;
``198 öğrenci tam iki saat çalışmıştır'' denemez. Sayısal 1--4 kodlarına
Scale atanması ölçümün anlamını değiştirmez. Okul tablosunda GP'nin 349,
MS'nin 46 kayıtla temsil edildiğini görmek, örneklemin hedef evreni temsil
ettiğini göstermez. \texttt{g3} ortalaması \SPContextMean\ bir örneklem
istatistiğidir; daha geniş bir evrenin ortalaması ayrıca tanımlanmalıdır.
\end{outputguidebox}


\par\noindent\begin{minipage}{\linewidth}
\subsection{R uygulaması}
\label{sec:r-gercek-b02}

\texttt{prop.table} frekansları orana çevirir. \texttt{ordered} çalışma süresi kodlarının sıralı olduğunu belirtir.

\begin{lstlisting}[style=prismR]
d <- read.csv("companion/spss/csv/b02.csv")
stopifnot(!anyNA(d))
d$school <- factor(d$school, 1:2, c("GP", "MS"))
d$sex <- factor(d$sex, 1:2, c("F", "M"))
d$studytime <- ordered(d$studytime, levels=1:4)
for (v in c("school", "sex", "studytime")) {
  f <- table(d[[v]])
  print(f); print(100 * prop.table(f))
}
print(sapply(d[c("age", "g3")], function(v)
  c(n=length(v), mean=mean(v), sd=sd(v),
    min=min(v), max=max(v))))
\end{lstlisting}

\textbf{Çıktıyı okuma.} Paylaşılan R çıktısı 395 satır ve 10 sütun
gösterir; her sütunda eksik değer sayısı sıfırdır. \texttt{n},
\texttt{mean}, \texttt{sd}, \texttt{min} ve \texttt{max} satırlarını
ortak tablonun sütunlarıyla eşleştirin. Frekanslar ve yüzdeler de
paylaşılan çıktıda doğrulanmıştır.
\end{minipage}\par

\par\noindent\begin{minipage}{\linewidth}
\subsection{Python uygulaması}
\label{sec:python-b02}

\texttt{normalize=True} kategori oranlarını verir. Okul ve cinsiyet kodlarının anlamı ortak veri sözlüğünden okunmalıdır.

\begin{lstlisting}[style=prismPythonApp]
import pandas as pd
import numpy as np
from scipy import stats

d = pd.read_csv("companion/spss/csv/b02.csv")
assert not d.isna().any().any()
for v in ["school", "sex", "studytime"]:
    print(d[v].value_counts().sort_index())
    print(100*d[v].value_counts(normalize=True).sort_index())
print(d[["age", "g3"]].agg(
    ["count", "mean", "std", "min", "max"]))
\end{lstlisting}

\textbf{Çıktıyı okuma.} Paylaşılan Python çıktısı da 395 satır,
10 sütun ve her sütunda sıfır eksik değer gösterir. \texttt{count},
\texttt{mean}, \texttt{std}, \texttt{min} ve \texttt{max} satırları
R özetleriyle eşleşir. Paylaşılan çıktıda okul ve cinsiyet kategorileri
etiketleriyle, yukarıdaki kısa kodda ise sayısal kodlarıyla gösterilir.
Bu betimsel uygulama bir hipotez testi veya $p$ değeri üretmez.
\end{minipage}\par

\subsection{Sonuçların karşılaştırılması ve ortak rapor}
12 Eylül 2026'da paylaşılan SPSS ekran görüntüsü ile R ve Python
metin çıktıları karşılaştırılmıştır. Üç kategorik değişkenin frekansları
eşleşmiş; yüzdeler gösterilen yuvarlama hassasiyetinde uyum sağlamıştır.
Yaş ve yıl sonu notuna ait 10 betimsel hücre R ve Python'da altı ondalık
basamakta eşleşmiş, SPSS'in yuvarlanmış değerleriyle de uyum göstermiştir.
Aşağıdaki tablolar bu çıktıların kitap düzenine uyarlanmış özetidir;
ekran görüntüsü değildir.

\begin{table}[htbp]
\centering
\caption{B02 verisinde okul, cinsiyet kodu ve çalışma süresi dağılımları}
\label{tab:b02-dogrulanmis-frekans}
\small
\begin{tabular}{@{}llrr@{}}
\toprule
Değişken & Kategori & Frekans & Yüzde \\
\midrule
Okul & GP & 349 & 88,4 \\
     & MS & 46 & 11,6 \\
\midrule
Cinsiyet kodu & F & 208 & 52,7 \\
              & M & 187 & 47,3 \\
\midrule
Çalışma süresi & 1: $<2$ saat & 105 & 26,6 \\
              & 2: 2--5 saat & 198 & 50,1 \\
              & 3: 5--10 saat & 65 & 16,5 \\
              & 4: $>10$ saat & 27 & 6,8 \\
\bottomrule
\end{tabular}
\par\smallskip
{\footnotesize Her değişken için toplam 395 kayıt ve yüzde 100'dür.
Yüzdelerin paydası 395'tir; eksik değer olmadığından SPSS'in
\emph{Percent} ve \emph{Valid Percent} sütunları eşittir.
Çalışma süresi etiketleri kaynak verideki sıralı kategorilerdir.\par}
\end{table}

\begin{table}[htbp]
\centering
\caption{B02 verisinde yaş ve yıl sonu notunun doğrulanmış özetleri}
\label{tab:b02-dogrulanmis-betimsel}
\small
\begin{tabular}{@{}lrrrrr@{}}
\toprule
Değişken & $n$ & Ortalama & Std. sapma & Min. & Maks. \\
\midrule
Yaş & 395 & 16,70 & 1,276 & 15 & 22 \\
\texttt{g3} & 395 & 10,42 & 4,581 & 0 & 20 \\
\bottomrule
\end{tabular}
\par\smallskip
{\footnotesize Ortalamalar iki, örneklem standart sapmaları üç ondalık
basamağa yuvarlanmıştır. SPSS'te ortak geçerli gözlem sayısı
(Valid N, listwise) 395'tir.\par}
\end{table}

\begin{outputguidebox}
Tablo~\ref{tab:b02-dogrulanmis-frekans} içinde en sık görülen çalışma
süresi kategorisi 198 kayıtla (\%50,1) haftalık 2--5 saattir.
\texttt{studytime} kodları saat miktarı olarak yorumlanmaz.
Tablo~\ref{tab:b02-dogrulanmis-betimsel} aynı öğrencilerin yaş ve not
dağılımlarını özetler. Üç yazılımın uyumu hesapların tutarlılığını
destekler; öğrencilerin rastgele seçildiğini veya Türkiye'deki lise
öğrencilerini temsil ettiğini kanıtlamaz. B01 ve B02 aynı kayıtları
kullandığından iki bağımsız örneklem gibi birleştirilmemelidir.
\end{outputguidebox}

\textbf{Raporlama örneği (doğrulanmış çıktılarla).}
``İki okuldan 395 öğrenci kaydının 198'inde (\%50,1) haftalık çalışma
süresi 2--5 saat kategorisindedir. Kategori kodları saat miktarı olarak
değil, sıralı sınıflar olarak değerlendirilmiştir. Yaş ortalaması
16,70 yıl ($s=1{,}276$), yıl sonu matematik notu ortalaması 10,42 puandır
($s=4{,}581$). Bu sonuçlar gözlenen kayıtları betimler; Türkiye'deki lise
öğrencilerinin tümüne doğrudan genellenmemiştir.''


\textbf{Görev.} Hedef evreni ``Türkiye'deki lise öğrencileri'' olarak seçen
bir araştırmacı için bu dosyanın neden yetersiz olduğunu açıklayın. Dosyayı
yalnız eğitim amaçlı sonlu bir evren kabul ettiğinizde parametrenin nasıl
farklı tanımlandığını belirtin.

\section{R ile çözümlü uygulama}
\label{sec:r-b02}

\textbf{Soru.} Bölümdeki dört öğrencilik örnekte okul türünü nominal,
sınıf düzeyini sıralı, puanı nicel değişken olarak tanımlayın. Devlet okulu
oranını ve ortalama puanı hesaplayın.

\begin{lstlisting}[style=prismR]
veri <- data.frame(
  okul_turu = factor(c("Devlet", "Ozel",
                      "Devlet", "Devlet")),
  sinif = ordered(c(1, 2, 1, 3), levels = 1:3),
  puan = c(72, 81, 68, 77)
)
str(veri)
prop.table(table(veri$okul_turu))
table(veri$sinif)
mean(veri$puan)
tapply(veri$puan, veri$okul_turu, mean)
\end{lstlisting}

\begin{outputguidebox}
\textbf{Çözüm ve kontrol değerleri.} Devlet okulu payı 0,75, özel okul
payı 0,25'tir. Sınıf düzeylerinin frekansları sırasıyla 2, 1 ve 1'dir.
Genel puan ortalaması 74,5; devlet okulundaki üç öğrencinin ortalaması
72,3333; özel okuldaki tek öğrencinin puanı 81'dir.
\end{outputguidebox}

\textbf{Yorum.} \texttt{ordered} sınıfların sırasını kaydeder; sınıf
kodlarının ortalaması başarı ortalaması yerine geçmez. Okul türüne göre
betimsel fark, tek başına okul türünün nedensel etkisi değildir.

\textbf{Pekiştirme ve yanıt.} Dört kayıt hedef evrenin tamamıysa 74,5
bir evren ortalamasıdır. Daha geniş bir evrenden seçilmişlerse aynı sayı
örneklem ortalamasıdır. Bu ayrımı R komutu değil araştırmanın kapsamı belirler.

\section{Bölüm özeti}

\begin{itemize}
  \item Hedef evren, erişilebilir evren ve örneklem farklı kümeler olabilir.
  \item Parametre evrene, istatistik örnekleme ait sayısal özettir.
  \item Kategorik ve nicel değişkenler farklı özet ve grafikler gerektirir.
  \item Sayısal kod, değişkenin ölçme düzeyini kendiliğinden değiştirmez.
  \item Örnekleme yöntemi ve örnekleme çerçevesi genelleme sınırını belirler.
  \item Büyük örneklem sistematik kapsam, gönüllülük veya yanıtlamama
  yanlılığını kendiliğinden gidermez.
\end{itemize}

\section*{Alıştırmalar}

\begin{enumerate}
  \item Hedef evren ile erişilebilir evrenin farklı olduğu bir araştırma örneği yazın.
  \item Parametre ve istatistiği aynı araştırma sorusu üzerinde gösterin.
  \item Nominal, sıralı, kesikli ve sürekli değişkenlere birer örnek verin.
  \item Sayısal kodlanmış her değişkenin neden nicel olmadığını açıklayın.
  \item Büyük örneklemin temsil gücünü tek başına neden garanti etmediğini tartışın.
  \item Hedef evren, erişilebilir evren ve örnekleme çerçevesini aynı eğitim
  araştırması üzerinde birbirinden ayırın.
  \item Öğrencilerin sınıflar içinde kümelendiği bir çalışmada gözlem
  bağımsızlığı varsayımının neden zayıflayabileceğini açıklayın.
  \item 250 öğrencinin 145'inin bir uygulamayı başarılı bulduğu örneklemde
  örneklem oranını hesaplayın ve parametreyi tanımlayın.
  \item Basit rastgele, tabakalı ve küme örneklemesini seçim birimi ve kullanım
  amacı bakımından karşılaştırın.
  \item Lisans, yüksek lisans ve doktora öğrencilerinin ayrı temsil edilmesi
  gereken bir araştırma için tabakalı örnekleme planı hazırlayın.
  \item Bir veri sözlüğünde bulunması gereken en az beş bilgi türünü yazın.
  \item Eksik değer için 99 kodunun kullanılmasının analizde doğurabileceği
  sorunu ve çözümünü açıklayın.
\end{enumerate}